% !TEX root = ../main.tex
%% Corrigé du sujet type bac n° 5 — valeurs vérifiées : verifications/sujets_bac.py

\corrigeSujet{5}

\begin{corrigeBac}{8}{La pile d'une aide auditive}
\partieExo{1}{Le fonctionnement d'une pile}
\begin{questions}
\item $Q_{r,i}=\dfrac{[\ce{Zn^{2+}}]_i}{[\ce{Cu^{2+}}]_i}=\dfrac{0{,}10}{0{,}10}=\res{1}$. Comme $Q_{r,i}<K$, le système évolue dans le \res{\text{sens direct}} : le zinc
est oxydé et les ions cuivre(II) sont réduits. \pts{1 pt}
\item Le zinc, oxydé, cède des électrons : la lame de zinc est la \res{\text{borne négative}}, la lame de cuivre la borne
positive. Dans le circuit extérieur, les électrons vont de la lame de zinc vers la lame de cuivre. Le pont salin ferme le
circuit et assure l'électroneutralité des solutions par la migration de ses ions. \pts{1 pt}
\item Au cours du fonctionnement, $Q_r$ augmente et se rapproche de $K$ ; la tension diminue. À l'équilibre, $Q_r=K$ : la pile
est usée et sa tension est \res{\text{nulle}}. \pts{0,5 pt}
\end{questions}
\partieExo{2}{La pile zinc-air}
\begin{questions}[resume]
\item Le zinc est oxydé : l'électrode de zinc est l'\res{\text{anode}}. Le dioxygène est réduit à la \res{\text{cathode}}, au contact de
l'air. En combinant les demi-équations (la première multipliée par 2) : $\res{\ce{2Zn + O2 -> 2ZnO}}$. \pts{1 pt}
\item $Q=0{,}620\times3600=\res{\qty{2232}{C}}$, soit $n(\mathrm{e}^-)=\dfrac{2232}{96500}=\qty{2,31e-2}{mol}$ ; $n(\ce{Zn})=\dfrac{n(\mathrm{e}^-)}{2}=\qty{1,16e-2}{mol}$, soit
$m=1{,}16\times10^{-2}\times65{,}4=\res{\qty{0,76}{g}}$. \pts{1,25 pt}
\item $\Delta t=\dfrac{620}{1{,}2}=\qty{517}{h}$, soit $\dfrac{517}{16}=\res{32}$ jours d'utilisation. \pts{0,75 pt}
\item $\mathcal{E}=U\,Q=1{,}4\times2232=\res{\qty{3,1e3}{J}}=\qty{0,87}{W.h}$, soit $\dfrac{0{,}87}{\num{1.9e-3}}=\res{\qty{4,6e2}{W.h.kg^{-1}}}$ : près de deux fois plus
qu'une batterie lithium-ion. Le réactif oxydant, le dioxygène, n'est pas stocké dans la pile : il vient de l'air. \pts{1 pt}
\item $n(\ce{O2})=\dfrac{n(\mathrm{e}^-)}{4}=\qty{5,8e-3}{mol}$, soit $5{,}8\times10^{-3}\times24=\qty{0,14}{L}$ de dioxygène, contenus dans $\dfrac{0{,}14}{0{,}21}=\res{\qty{0,66}{L}}$
d'air. Tant que les trous sont obturés, le dioxygène ne peut pas entrer : la pile ne fonctionne pas et ne se décharge
pas pendant son stockage. Une fois la pastille retirée, la réaction démarre, même si l'appareil est éteint. \pts{1,5 pt}
\end{questions}
\end{corrigeBac}

\begin{corrigeBac}{7}{Le saut en parachute}
\begin{questions}
\item Le poids $\vec P$, vertical vers le bas, de valeur $m\,g$, et la force de frottement $\vec f$, verticale vers le haut, de valeur
$k\,v^2$. \pts{0,5 pt}
\item Deuxième loi de Newton projetée sur l'axe vertical orienté vers le bas : $m\,\dfrac{\mathrm{d}v}{\mathrm{d}t}=m\,g-k\,v^2$, d'où
$\res{\dfrac{\mathrm{d}v}{\mathrm{d}t}=g-\dfrac{k}{m}\,v^2}$. \pts{1 pt}
\item Lorsque la vitesse limite est atteinte, $\dfrac{\mathrm{d}v}{\mathrm{d}t}=0$ : $v_\mathrm{lim}=\sqrt{\dfrac{m\,g}{k}}=\sqrt{\dfrac{80\times9{,}81}{0{,}25}}=\res{\qty{56}{m.s^{-1}}}$, soit
$\res{\qty{202}{km.h^{-1}}}$. \pts{1 pt}
\item Méthode d'Euler : sur un petit intervalle $\mathrm{d}t$, $v(t+\mathrm{d}t)\approx v(t)+\dfrac{\mathrm{d}v}{\mathrm{d}t}\,\mathrm{d}t$, avec $\dfrac{\mathrm{d}v}{\mathrm{d}t}=g-\dfrac{k}{m}\,v^2$. La boucle
s'arrête quand $v$ atteint 95~\% de $v_\mathrm{lim}$ : il faut \res{\qty{10,4}{s}}, la vitesse vaut alors \qty{53,3}{m.s^{-1}} et le parachutiste a
parcouru \res{\qty{373}{m}}. \pts{1,25 pt}
\item La tangente à l'origine passe par environ $(\qty{5}{s}~;~\qty{49}{m.s^{-1}})$ : sa pente vaut \qty{9,8}{m.s^{-2}}. À $t=0$, la vitesse est
nulle, donc les frottements aussi : seul le poids agit, et $a=g$. \pts{0,5 pt}
\item Sans frottement : $z=\frac12\,g\,t^2=\frac12\times9{,}81\times10{,}4^2=\res{\qty{531}{m}}$, au lieu de \qty{373}{m}. \pts{0,75 pt}
\item L'énergie potentielle perdue est transférée à l'air et au parachutiste sous forme thermique par les frottements.
Puissance : $f\,v=m\,g\,v_\mathrm{lim}=80\times9{,}81\times56=\res{\qty{4,4e4}{W}}$. \pts{0,75 pt}
\item $v'_\mathrm{lim}=\sqrt{\dfrac{80\times9{,}81}{30}}=\res{\qty{5,1}{m.s^{-1}}}$. Une chute libre sans frottement depuis la hauteur $H$ donne
$v=\sqrt{2gH}$, d'où $H=\dfrac{v^2}{2g}=\dfrac{5{,}1^2}{2\times9{,}81}=\res{\qty{1,3}{m}}$ : atterrir équivaut à sauter d'un muret. \pts{1,25 pt}
\end{questions}
\schema[corrige]{bac-s5-parachute}
\end{corrigeBac}

\begin{corrigeBac}{5}{Le diamètre d'un cheveu}
\begin{questions}
\item La diffraction est d'autant plus marquée que l'obstacle est petit : avec un cheveu, l'angle vaut environ
\qty{0,01}{rad}, ce qui donne sur un écran éloigné une tache de quelques centimètres. Le laser, monochromatique et
directif, rend la figure nette. \pts{0,5 pt}
\item Dans le triangle formé par le cheveu et un bord de la tache : $\tan\theta=\dfrac{L/2}{D}\approx\theta=\dfrac{\lambda}{a}$, d'où
$\res{a=\dfrac{2\,\lambda\,D}{L}}$. \pts{1 pt}
\item $a=\dfrac{2\times\num{650e-9}\times2{,}00}{0{,}037}=\qty{7,0e-5}{m}$, soit $\res{\qty{70}{\micro m}}$. \pts{0,75 pt}
\item $\dfrac{u(a)}{a}=\sqrt{\left(\dfrac{0{,}1}{3{,}7}\right)^2+\left(\dfrac{0{,}01}{2{,}00}\right)^2}=2{,}7~\%$ : $u(a)=\qty{1,9}{\micro m}$ et $\res{a=\qty{70,3}{\micro m}}$ avec
$\res{u(a)=\qty{1,9}{\micro m}}$. \pts{1 pt}
\item L'écart vaut $72-70{,}3=\qty{1,7}{\micro m}$, d'incertitude-type $\sqrt{1{,}9^2+1^2}=\qty{2,1}{\micro m}$ : $z=\dfrac{1{,}7}{2{,}1}=0{,}8<2$. Les deux
mesures sont \res{\text{compatibles}}. \pts{1 pt}
\item Éloigner l'écran, pour agrandir $L$ et réduire son incertitude relative ; mesurer la distance entre des minimums plus
éloignés ou photographier la figure et l'exploiter avec un logiciel ; répéter la mesure pour une évaluation de type A. \pts{0,75 pt}
\end{questions}
\end{corrigeBac}
